\subsection{An isolated stationary neutral band}

Consider the periodic eight-cell fermion, quantum-scalar and 119-state-link
source with nonblank-link coefficient two and hopping coefficient
$\lambda=2^{-34}$. Every local scalar and Yukawa coefficient is retained.
In the physical sector with raw lepton number seven and baryon number zero,
the atomic band has one local neutral singlet of energy $e_\nu$ and
seven local ground doublets of energy $e_0$. It has rank 1024, with rank 128 at each of eight momenta.
Its complement gap exceeds $d=0.0110299637715905$ and the hopping
perturbation has norm at most $v=360\lambda$. Blank-link selection gives
$PTP=0$. A Schur-complement bound therefore confines the perturbed band to
\[
 [a-\epsilon,a],\qquad a=7e_0+e_\nu,\qquad
 \epsilon=\frac{v^2}{d-v}\leq3.9810\,10^{-14}.
\]
The spectral projector satisfies $\|E-P\|\leq v/d<1.900\,10^{-6}$.
The same source's positive second-order exchange and bounded quartic
remainder prove a unique physical ground at zero momentum; the
certified ground gap exceeds $4.4542\,10^{-25}$.
Relative to that ground, the neutral band is contained in
$[0.0895553304760346,0.0895553304760794]$.
This argument uses the actual finite exchange operator and a uniform
remainder, avoiding diagonalization of the entire interacting carrier.

The local neutral rank-one projector is an explicit algebraic polynomial
of the sixteen-dimensional native local block and is annihilated by all
twelve gauge currents. On the full carrier, a Chebyshev filter with
repeated rejection gives a degree-167608 polynomial $M(H)$ approximating
the neutral band projector, with operator error below $10^{-69}$.
Together with the commuting gauge, charge and momentum projectors, it
defines the complete accept/fail instrument. From the maximally mixed
input, the accepted source density commutes exactly with $H$ and is
within $7.02\,10^{-13}$ in unhalved trace norm of the uniform
zero-momentum band. The original-input success probability is approximately
$2.85\,10^{-126}$. The polynomial degree counts Hamiltonian evaluations;
it is not a primitive gate count or a claim of full-matrix execution.
Stationarity concerns the reduced source density and its time-independent
source readouts, without asserting stationarity of an arbitrary purification.

A complete 51-record instrument retains preparation failure, energy and
profile complements, and 48 position/occupation outcomes. The stationary
conditional profile is uniform. A separate localized filtered preparation
has leakage bounded, up to the specified filter error, by
\[
 \left(\frac{\eta}{\sqrt{1-\eta^2}}+\epsilon |t|\right)^2,
 \qquad \eta=v/d.
\]
At $t=10^{12}$ this bound is below $0.001585$. Position and momentum
projectors do not commute; their probabilities are separate instruments.
The entire band is almost flat, so these bounds supply no positive
propagation speed. A unit-speed relativistic assignment is incompatible
with its momentum width in the declared source units.

Calibrating a separate native charged local gap by the empirical electron
energy yields $6.61638659224\,10^{-5}\,\mathrm{GeV}$ for this neutral
band. It is a calibrated zero-momentum excitation energy. Neither a
relativistic mass shell nor a dark-matter or neutrino identification is
inferred from its stability and neutrality.

\subsection{Self-consistent magnetic precession and refinement}

For the confined positive kernel $K(A)=|D_W(A)|$, use the semiclassical
energy $\langle\psi,K(A+A_{\rm ext})\psi\rangle+
 a^3\|CA\|^2/(2\kappa)$ on mean-zero transverse link fields.
The full absolute-kernel derivative retains both signs of the Wilson
operator. In its eigenbasis the density entering the current is
\[
 W_{ij}=\frac{\lambda_i+\lambda_j}{|\lambda_i|+|\lambda_j|}\rho_{ij},
 \qquad J_e=-\operatorname{Tr}(W\partial_eD_W).
\]
Only the transverse projection of $J$ sources the quasistatic magnetic
field; the full current obeys the source continuity law. A uniform
Hessian estimate on the executed coarse cube makes the fixed-state
field map a contraction with factor below $0.608135$. It proves a
unique field minimizer for each normalized matter state, without
asserting contraction of a coupled time-stepping algorithm.

Evolution under this same kernel gives a complete weak-field precession
cycle. The spin and self-sourced dipole frequencies are approximately
$0.002848096149$ and $0.002848053578$, while an independently optimized
stationary energy split is $0.002848114846$. Reversed, zero, halved and
frozen-field probes and step refinement retain their respective errors.
Complete energy, spin and moment instruments include all spectral
branches. These dynamics are nonlinear quasistatic mean-field dynamics;
they do not define a quantum channel on arbitrary entangled inputs.
Their magnetic field has no propagating electromagnetic degrees of freedom.

At fixed matter extent six and field extent twenty-one, spacings
$1,3/5,3/7,1/3$ give static bound-energy factors approximately
$1.47517,1.78875,1.88532,1.92738$. Independent volume and symmetric-probe
controls are retained. Fit-order-dependent intercepts ranging from
$2.0002$ to $2.0746$ do not give a certified continuum uncertainty.
On the finest matter grid a separately declared coupling
$\kappa=10^{-14}$ makes the self-field error uniformly small. A zero-probe stationary current witness on this same weak-coupling member
has eigen-residual below $7.5\,10^{-15}$ and a nonzero self-field. Its
static doublet normalization $1.92095809$ agrees with the weak-field
precession value to the stated finite-probe precision. Two
actual full-operator eigenline witnesses prepare the precessing state.
Exact rational link-phase Taylor bounds and tridiagonal Sturm counts
certify their residuals and isolate the complete low quartet.
Duhamel's identity bounds the full nonlinear trajectory through $T=6000$,
including every unprepared mode. The finite-duration ground-rest-energy
normalization lies in
\[
 1.920947397728\leq g_T\leq1.920952163162.
\]
The static convention divides by the measured canonical spin; the
precession convention uses its rotation frequency. Their numerical
values must not be equated when the projected spin is below one half.
The regulator, confinement and weak-coupling choice are supplied.
These finite results do not identify the precision physical electron
magnetic moment.

\subsection{Exact pinned Coulomb minima}

Replace the unit-link penalty in the periodic gauge source by
$(\kappa/2)\rho^Y G\rho^Y$, where $G$ is the mean-zero inverse of the
same lattice incidence Laplacian. Pin a complete native hypercharge
$-1$ weak/colour-singlet multiplicity space at $y$, a unit reference
at zero, and complete gauge-singlet matter multiplicities elsewhere.
Every link state and local multiplicity is retained. Each native hopping
word changes a pinned local density, so its compression vanishes exactly.
The constrained physical carrier factors into local multiplicities and
a nonempty flux space. Actual native strings prove nonemptiness and
saturate the tensor lower bound. Thus
\[
 E_{\min}(y)=(V-1)E_0+e_\chi+4+
                       \kappa(G_{00}-G_{0y}).
\]
The unknown complete charged local minimum $e_\chi$ cancels from every
separation difference. This removes the need to solve that local spectrum
in order to calculate the static force. Flux-loop degeneracy and the
full screening complement remain present; the pinning apparatus is
explicit and its carrier is not invariant under unrestricted hopping.

Finite-box Green calculations give the attractive Coulomb force. On a
side-128 box the force ratio to its continuum form is approximately
$1.00898$ at separation twelve; a potential-difference ratio is
$1.00407$ at separation eight. The supplied $\kappa$, physical length
conversion and original source's different link energy are distinguished.
No autonomous pair motion or empirical fine-structure constant follows
from this constrained minimum identity.

\subsection{Kinetic evolution in the coherent lepton compiler}

The coherent phase cycle extends to a homogeneous periodic kinetic law
without switching that law off. Let $K$ have the native nearest-neighbour
Weyl symbol $-\boldsymbol\sigma\cdot\sin(a\mathbf p)/a$ on every
internal, character and address copy. Ancillary-number phase rotations
$R=\exp(i\pi N_A/2)$ commute with $K$, so the four conjugated generators
$H_r=K+R^{-r}VR^r$ average to $K+\overline V$. The ancillary-vacuum
embedding therefore obeys
\[
 (K+\overline V)J=J(K_{\rm physical}+H_Y)
\]
on the complete physical Fock carrier, including external quantum
references. For each character, the physical address and its image under
the rank-one analyzer span an invariant two-dimensional space. The
noncommuting calculation acts on $24s$ fermionic modes at $s$ sites;
every other enlarged mode evolves as an exact free kinetic spectator.
This is an invariant-factor identity, rather than a discarded subspace.

Spatial inversion $I:x\mapsto-x$ reverses $K$ and preserves uniform
onsite pairing. Parity $S$ of the complete charged-lepton doublet leg
reverses $V$ and preserves $K$. The exterior lift of the site permutation
includes every fermionic occupation sign. Thus $C=IS$ satisfies
$CH_rC^*=-H_r$. This implements negative full-generator segments through
declared coherent permutation and parity services; leg parity alone
cannot do so. Exact finite Fock tests include that failing control.

On the invariant active factor, $\|H_r\|\leq B=36s/a+6s$.
The same fourth-order word identity applies with
$R_0=7B/3$ and gives complete unhalved output trace error
\[
 \epsilon\leq\min\{2,R_0^5|t|^5/(10n^4)\},\qquad
 n\geq R_0|t|.
\]
The free spectators accumulate exactly the same signed target time in
both evolutions. The controls and their time/work costs remain apparatus
inputs. The result retains the register's three coefficients and its
phase, scalar-background and calibration ancestry; it does not select
an interacting Higgs vacuum or eliminate massless spectators.

\subsection{Charged branch masses and complete spatial records}

The analyzer-derived pairing admits three explicit intertwiners from
the complete native Nambu source to charged Dirac frames. The source
Hamiltonian, electric current and all three spin generators restrict
respectively to $h_k=-\tau_3\otimes\sigma\cdot\pi+e_k\tau_1\otimes I$,
$Q=-I$ and $S_j=I\otimes\sigma_j/2$. Thus
$E_k^2=|\pi|^2+e_k^2$, with strict branch minima $e_k$ at each lattice
valley. Complete reservoir rest records and spatial positive-band probe
records evolve under that same quadratic Hamiltonian; time-zero
comparisons separate motion from the initial projected tail.

The side-four realization retains all $5760$ physical CAR modes and
eight valleys. Its three distinct positive masses therefore describe
$24$ family/valley copies, with $39$ additional massless internal
coefficients. An explicit three-quark colour-singlet zero-mode state has
unit electric charge and zero energy. It rules out a positive global
charged-sector mass floor in this charged-lepton-only action. A single
electron energy anchor converts the other two branch values to
approximately $0.1056594144959$ and $1.7769849708134\,\mathrm{GeV}$;
the retrospective phase choice, interacting particle identities and
physical clock are separate scientific conditions.

\subsection{Native spin, field and frame return periods}

On the periodic native kinetic and pairing construction, a literal spatial
quarter-turn combined with its spin rotation preserves every Hamiltonian
pulse. Its fourth power is fermion parity on the full carrier, including
analyzer coordinates. An odd charged packet therefore acquires an
interferometrically measurable minus sign at $2\pi$ and returns at $4\pi$,
relative to an idle pointer and charge-compensating reference. The filled
negative-band reference is invariant under proper cube rotations. Rotating
the opposite reference too cancels the sign, and pointer dephasing removes
the interference.

For an axial positive-band packet, the native spin current on its origin
plane is $s=(p/2)(\cos\theta,\sin\theta,0)$ under actual cube rotations,
with $p>0$ at every source time. An oriented frame constructed from this
current has period $2\pi$. The apolar tensor
$Q_{ij}=s_i s_j-\delta_{ij}|s|^2/3$, measured using two independent source
copies, has period $\pi$. The complete classical tensor law includes all
preparation failures and detector complements; its shorter period is not a
period of the complete quantum source state. For a spin-up axial preparation,
the same native excess density integrates to $1/2$ throughout autonomous
evolution. The unfiltered occupied sea retains a nonzero spin variance,
so this integral does not imply a total-spin eigenstate or a continuous
spatial Noether angular momentum. The field is a declared two-copy observable,
and the rotation theorem concerns exact cube symmetries of the finite source.

\subsection{Fractional charge with native colour strings}

The complete particle source also supports a direct quark-family
charge/flux experiment. The specified preparation fills all 240 source
fermion modes, leaves the scalar and link registers blank, and prepares
the thirty-mode reference vacuum. It is a supplied highly occupied
gauge singlet, with no assertion that it is the interacting ground.
For either weak-singlet colour triplet, the reference-hole pump
\[
 B=\frac1{\sqrt3}\sum_i a_{{\rm ref},i}^*c_{0,i}
\]
has $B^*B=N_{\rm family,spin\,0}/3\leq I$ on reference-vacuum
inputs and succeeds with probability one on this preparation. Its full
failure complement is retained.

Contracting the actual 119-state link coordinates along a path of
$\ell$ edges gives normalized coefficients $3^{-(\ell+1)/2}$.
The endpoint is the literal CAR vector $c_i|\mathrm{filled}\rangle$,
including its occupation sign. The reference, link and dual endpoint
generators cancel at every vertex. The source-hole charges are
$+2/3$ and $-1/3$; their weak charge is zero and their colour flux
continues along the string. The calculation checks all $6{,}648$
shortest-path constraints and $1{,}536$ constraints using the original
full CAR/scalar/link Gauss compiler.

For the native colour Cartan $C$ induced by
$\operatorname{diag}(1,1,-2)/3$, the complete Fock identity is
$e^{2\pi i d\Gamma(Q)}e^{2\pi i d\Gamma(C)}=I$.
Both selected holes have electromagnetic two-pi phase
$-1/2-i\sqrt3/2$, completed by the inverse colour-centre phase.
A controlled endpoint hypercharge rotation measures this phase without
choosing a gauge-variant colour component. Its complete randomly chosen
pointer quadratures have probabilities
\[
 \left(\frac18,\frac38,\frac{2-\sqrt3}{8},
                         \frac{2+\sqrt3}{8}\right).
\]
The three-turn electromagnetic control gives
$(1/2,0,1/4,1/4)$. This is an internal gauge-path holonomy; the
electromagnetic path does not close in one turn and supplies no
fractional spatial texture winding.

The original full hopping operator places squared amplitude $1/16$
on each of six one-edge arrival subspaces, retaining both endpoint
spins, and has total squared norm $3/8$ on the prepared vector.
Each arrival projector annihilates the no-hopping evolution because
the atomic Hamiltonian preserves blank links. With
$\|H_0\|\leq480$, $\|T\|\leq360$, $\lambda=2^{-34}$ and
$t=2^{-24}$, Duhamel's identity encloses each original-input arrival
probability in approximately
$[6.91604\,10^{-37},8.15583\,10^{-37}]$.
The full complement, reference and pump failures are retained. The
positive finite-time signal does not identify a low-energy quark,
its rest mass, a hadron or an unrestricted confinement phase.

\subsection{A measured scale-dependent colour coupling}

The colour restriction of the five-level seam cost is exactly
\[
 1-\tfrac15\operatorname{Re}\operatorname{tr}_5
       \operatorname{diag}(U,I_2)
 =\tfrac35\left(1-\tfrac13\operatorname{Re}\operatorname{tr}_3U\right).
\]
A declared four-dimensional periodic thermal construction with cost
multiplier ten therefore gives Wilson weight $\beta=6$, or bare
$g_0^2=1$ in the fundamental trace convention. The native generators
give $\operatorname{tr}(T_aT_b)=\delta_{ab}/2$ and $C_F=4/3$.
This colour ensemble and its thermal state are declared separately
from the finite particle Hamiltonian; no observed strong coupling is
an input.

Let $W(R,T)$ be the ensemble mean of the normalized real trace of
a rectangular Wilson loop, averaged over all sites and coordinate
plane orientations. Define
\[
 \chi_R=-\log\frac{W(R,R)W(R-1,R-1)}{W(R,R-1)^2},\qquad
 g_C^2(R;N)=\frac{\chi_R}{k_R(N)}.
\]
The coefficient $k_R$ removes the free finite-box geometry. If $J_C$
is the conserved oriented link current of the rectangle and
$\widehat k^2=4\sum_\mu\sin^2(k_\mu/2)$, the exact finite Gaussian
covariance gives
\[
 f(R,T)=\frac{C_F}{2N^4}\sum_{k\ne0}
              \frac{\sum_\mu|J_{C,\mu}(k)|^2}{\widehat k^2},
 \qquad
 k_R=f(R,R)+f(R-1,R-1)-2f(R,R-1).
\]
In particular $k_1=(1-N^{-4})/3$. An independent position-space inverse
verifies the coefficient. The Gaussian control $W=e^{-g^2 f}$ gives
the same $g_C^2=g^2$ at every scale; the unnormalized loop ratio would
produce a geometric false trend. This defines the normalization scheme
without assuming an interacting finite-volume zero-mode expansion.

Four cold/hot productions at the same $\beta=6$ retain 400 records
each on sides six and eight. Common-block deletions preserve loop and
scale correlations. At side eight, the combined coupling squared is
approximately $1.56186$, $2.73542$, $4.13672$ at lengths
$0.5$, $1.5$, $2.5$ in lattice units; twenty-record block errors are
$0.00050$, $0.00444$, $0.03837$. Side-six values are
$1.55976$, $2.75157$, $5.43421$. The longest-loop discrepancy exposes
substantial finite-volume sensitivity. Multiple block sizes, independent
starts, split histories, all raw records and gauge controls are retained.
These errors are sampling diagnostics, not rigorous mixing or confidence
bounds. No imported perturbative running curve enters the calculation.
The finite coupling weakens toward shorter measured lengths; a continuum
beta function, physical energy scale and onset at a particle core are
not identified.
